\begin{abstract}
Selecting observations with correlated errors changes the covariance that must
be inverted. We develop global certificates for the resulting discrete
Gaussian information-design problem, with known covariance and nominal mean
sensitivities. Genuine local Gaussian conditionals yield explicit relative
precision bounds, uniform over selected schedules, for noisy scalar chains,
complete observed blocks, general exponentially decaying covariances, and
fixed partial-observation packets. The bounds transfer finite-state support
prices to the original design objective and give a rational weighted-trace
approximation scheme under fixed decay promises. For fixed information
dimension and additive rational positive-semidefinite atoms, we also construct
polynomial sets of actual feasible paths in explicit acyclic graphs or
rationally represented matroid bases that approximate every attainable
information matrix in two-sided relative positive-semidefinite order,
preserving singular ranges exactly. This supplies determinant, minimum
eigenvalue, inverse-trace, and estimable-contrast consequences. For practical
log-determinant bounds, we provide rational witnesses against every admissible
diagonal virtual-noise split, finite-scenario certificates with uncertain
optimal normalizers, and a nested latent-anchor hierarchy proved by Schur
elimination. On ten stipulated rational instances, finite-history certificates
are strictly stronger than every scalar-split relaxation; one kinetic instance
also separates all diagonal splits. Independent enumeration of a public
kinetics model checks 2347 schedules and identifies one changed trace-optimal
choice under the correct information formula. Exact small studies verify
strict anchor-hierarchy improvements, while full-packet and fixed-physical-grid
experiments show limits on practical strength and cost. A standalone supplement
replays the certificates without a commercial solver.
\end{abstract}

\noindent\textbf{Keywords:} optimal experimental design; correlated errors;
discrete optimization; relative matrix approximation; global certificates;
latent Gaussian models.

\section{Introduction}\label{sec:introduction}

Measurement selection combines a statistical decision with a discrete resource
allocation problem. A kinetic experiment may choose which species to measure,
when to sample, and which instruments to install. A sensor schedule may acquire
an entire multichannel packet whenever a device is activated. Installation
costs, exact sample counts, and minimum time gaps restrict the feasible
experiments. Correlation couples the value of these choices: the information
from a new observation depends on which other observations are acquired.
Correlated dynamic multiresponse design has a long history
\citep{Patan2007}, and recent process-design formulations integrate measurement
costs with model-based information criteria \citep{Wang2024}.
For optimization, the central question is how to bound the best attainable
information over all feasible discrete experiments.

Write $F_S$ for nominal mean sensitivities at the selected observations,
$R_{SS}$ for their known positive definite error covariance, and $J_0$ for
fixed prior information or a declared regularizer. The information matrix is
$J(S)=J_0+F_S^TR_{SS}^{-1}F_S$. The selection is made before observing the
data, and it changes neither the process nor its covariance. Even with a
linear mean this matrix is generally not a sum of fixed per-observation
contributions. Extracting selected entries of the full precision $R^{-1}$
instead gives a different information model, as already explained by
\citet{Liu2016}. Section~\ref{sec:foundations} establishes the exact comparison
and examines its consequences in a pinned public implementation. The issue
is relevant to interpreting a solver's gap: a valid optimization bound must
refer to the statistical experiment that will actually be performed.

There is a second distinction even when the information formula is correct.
A continuous mixture of feasible schedules can have a larger log determinant
than any one schedule. Closing its numerical optimization gap does not close
this integrality gap. Approximating the covariance introduces a further error
that must be controlled uniformly over the choices made by the optimizer.
We keep these effects separate. A certificate consists of a feasible design
with a checked objective lower bound and a global upper bound for the same
discrete model. For $p$ parameters, a log-determinant gap $g$ certifies
determinant-root efficiency at least $e^{-g/p}$. This interpretation is useful
whether the design came from a heuristic, a convex relaxation, or an exact
combinatorial search.

\subsection{Relationship to prior work}\label{subsec:intro-prior}

Three established lines of research provide the starting point.
First, virtual-noise design and covariance-split sensor relaxations supply
concave extensions of selected-covariance information
\citep{Liu2016,Pazman2022}. \citet[Proposition 3]{Hainy2025} prove their
equivalence and develop practical optimization by simplicial decomposition.
General information-matrix hulls and subsystem criteria already admit support
conditions and conic formulations \citep{HarmanTrnovska2009,Sagnol2015}.
These results make both credible comparators and useful building blocks.
Our comparisons concern the strength of specified relaxation families, rather
than treating all dense formulations as one algorithm.

Second, local conditional Gaussian factorizations originate with
\citet{Vecchia1988}. Their noisy and vector extensions, sparse inverse-factor
interpretations, and approximation theory are substantial existing work
\citep{KatzfussGuinness2021,Schafer2021,Huan2025,KaminetzWebber2026}.
For a fully observed error process that is itself Markov, exact gap
conditionals lead to an ordered path representation. Its inverse-matrix hull
is covered by the convexification results of \citet{Lee2026}, building on
\citet{Wei2023}. Independent measurement noise generally removes this observed
Markov property. The finite-history analysis here addresses that noisy case
through the actual local regressions, with explicit constants for chronological
calendar windows. It does not introduce a new factorization or a general
localization principle.

Third, bounded-width Gaussian message schemes, approximate Pareto sets,
nonlinear matroid optimization, and fixed-dimensional experimental-design
approximation provide algorithmic precedents
\citep{Mahalanabis2012,Papadimitriou2000,Tsaggouris2009,Berstein2008,Brown2024}.
Their representations and quantifiers matter. A scalar focused-target
reduction already yields an approximation scheme from the Gaussian-message
predecessor; Appendix~\ref{app:scalar-prior} gives the required adaptation.
The spectral construction in this paper instead approximates each feasible
matrix by one actual feasible object in both directions of positive-semidefinite
(PSD) order. Its
matroid theorem requires an explicit rational representation and additive
matrix atoms. It is not a result for arbitrary correlated matroid selection.

Table~\ref{tab:prior-contributions} identifies what is inherited and what is
developed here. Detailed comparisons and counterexamples accompany the relevant
proofs, including the distinction between total Gaussian KL error and relative
precision error, and the importance of exact singular ranges. Literature
comparisons use the inspected primary versions identified in the bibliography;
unavailable final texts are not used to infer the absence of a result.

\begin{table}[htbp]
\centering\small
\caption{Contribution boundaries. The new results combine established
statistical and optimization machinery under the explicit hypotheses given
in the cited sections.}
\label{tab:prior-contributions}
\begin{tabular}{@{}>{\raggedright\arraybackslash}p{.20\textwidth}>{\raggedright\arraybackslash}p{.34\textwidth}>{\raggedright\arraybackslash}p{.40\textwidth}@{}}
\toprule
Component & Established antecedent & Development in this paper\\
\midrule
Selected information & Selected covariance and gating distinction
\citep{Liu2016}; classical Kantorovich inequality \citep{Moradi2019}
& Sharp design-efficiency consequences, equality/rank analysis, and exact
public-model reanalysis (Sections~\ref{sec:foundations},
\ref{subsec:source-computation}).\\[4pt]
Finite history & Vecchia conditionals, inverse localization and covariance
transport \citep{Vecchia1988,Schafer2021,Kozdoba2019}
& Explicit schedule-uniform relative bounds for four covariance classes,
spacing refinements, and transfer to exact feasible support prices
(Sections~\ref{sec:locality}, \ref{subsec:local-certificates}).\\[4pt]
Spectral approximation & Pareto/profile algorithms and guessed-vector
normalization \citep{Berstein2008,Tsaggouris2009,Brown2024}
& Deterministic all-target, two-sided PSD covers with polynomial accuracy
dependence and exact ranges for explicit rational acyclic graphs and represented matroids
(Section~\ref{sec:approximation}).\\[4pt]
Global comparison & Virtual-noise concavity and equivalence
\citep{Liu2016,Pazman2022,Hainy2025}; inverse-logdet convexity
\citep{KimKim2006}
& Finite rational witnesses separating every scalar or diagonal split,
including their stated affine-support intersections
(Section~\ref{subsec:all-split}).\\[4pt]
Robustness and anchors & Efficiency normalization
\citep{Burclova2016,Maus2010}; nuisance and subsystem design
\citep{Alexanderian2021,Sagnol2015}
& Certified uncertainty in optimal normalizers for one common schedule, and a proved
ordering across nested changing-anchor representations
(Sections~\ref{subsec:robust-certificate}, \ref{subsec:separators}).\\
\bottomrule
\end{tabular}
\end{table}

\subsection{Results and their assumptions}\label{subsec:intro-results}

The main theoretical link is a relative matrix sandwich. We prove explicit
bounds on the normalized residual covariance of fresh local conditionals,
simultaneously for every selected schedule. For noisy scalar and complete-block
chains, the constants use contraction and signal-to-noise bounds; the block
bound allows noncommuting transitions. General covariance decay needs an
off-diagonal exponential bound and a positive spectral floor. Fixed partial
observations admit a separate covariance-contraction argument, including
singular latent covariances. Theorems~\ref{thm:scalar-locality},
\ref{thm:block-locality}, \ref{thm:general-decay}, and
\ref{thm:partial-locality} give the constants. Combined with a calendar-state
graph, they yield a rational weighted-trace fully polynomial-time approximation
scheme (FPTAS) with input-sized observation
and parameter dimensions under fixed promises
(Theorem~\ref{thm:trace-fptas}). Exact count, mandatory/forbidden choices and
minimum spacing are included. Selecting channels independently within a
packet has a different complexity boundary
(Proposition~\ref{prop:channel-hardness}).

For fixed information dimension, Theorems~\ref{thm:dag-spectral-set} and
\ref{thm:matroid-spectral-set} return polynomial sets of actual paths or bases.
For every feasible target matrix, one returned matrix lies between its
$1-\eta$ and $1+\eta$ multiples in PSD order, even when it is singular.
The proof combines rational normalization on exact ranges with bounded signed
profiles and exact feasible-witness recovery. To the best of our knowledge,
this combined guarantee, with deterministic polynomial dependence on
$1/\eta$ under the stated explicit representations, is not stated in the
inspected predecessors. The result covers D, E and conventional A criteria, with consequences
for estimable contrasts; it requires no condition-number bound on the
information matrices. The polynomial degree grows with information dimension,
and the computational study does not implement these large spectral sets.

The practical certificates use a different consequence of the same matrix
control: an arbitrary positive definite reference and a globally priced
linear support give an upper bound for the original discrete objective
(Theorem~\ref{thm:local-certificate}). Rational arithmetic checks the proposed
reference, the full feasible support, and outward logarithm bounds. We extend
this construction to a shared finite-scenario schedule while enclosing each
scenario's unknown optimal normalizer
(Proposition~\ref{prop:robust-normalizers}). Separately, a fixed fractional
point and a repaired PSD dual give a lower witness for every admissible
diagonal virtual-noise relaxation (Proposition~\ref{prop:diagonal-barrier}).
A strict separation from a valid integer upper bound therefore survives both
complete continuous optimization and split tuning within that family.

When sufficient calendar histories are expensive, latent anchors provide an
exact conditional-block representation. Removing anchors gives a tighter
complete-schedule hull, as proved by the cross-representation Schur identity
and Theorem~\ref{thm:separator-hierarchy}. To the best of our knowledge, this
particular nested Markov-anchor hierarchy is not stated in the inspected
predecessors. The novelty claimed is this comparison between changing latent
representations and its explicit support certificates, within established
Schur and subsystem-design theory. Larger blocks cost more to enumerate, and
even the empty-anchor information hull can retain an integrality gap.

The numerical evidence in Section~\ref{sec:computation} tests these distinctions
on common input models. Ten exact comparisons separate every scalar split;
one kinetic comparison also separates all diagonal splits. New exhaustive
experiments verify strict nested-anchor improvements and a stronger block
locality constant. Public-source enumeration, robust standardization, partial
observation, and fixed-physical-grid studies supply complementary evidence.
The unsuccessful probes and cases favoring dense bounds are retained.
Section~\ref{sec:discussion} interprets their practical scope.
Throughout, covariance is known and parameter independent, sensitivities are
local, and a certificate concerns the supplied finite experiment. PSD priors
are allowed in the locality and spectral theory; the practical logdet and
separator certificates assume a positive definite prior. These guarantees do
not imply global nonlinear identifiability or robustness outside a declared
scenario set.
